% ============================================================
% PRN232 - Section 6: Database Design (Vietnamese)
% ============================================================
\section{\bi{Thiết kế cơ sở dữ liệu}{Database Design}}
\label{sec:db}

\subsection{\bi{Cơ sở dữ liệu quan hệ với Prisma}{Relational Database with Prisma}}

\bi{
Toàn bộ dữ liệu nghiệp vụ nằm trong PostgreSQL, được định nghĩa bằng
21 model Prisma trong \texttt{backend/prisma/schema.prisma}. Prisma
đóng vai trò lớp trừu tượng hoá: mọi truy vấn được viết bằng API
có kiểu \texttt{findMany}, \texttt{upsert} hoặc \texttt{createMany} và
được biên dịch thành SQL tham số hoá, nhờ vậy loại bỏ hoàn toàn khả
năng chèn SQL trực tiếp từ mã ứng dụng.
}{
All business data resides in PostgreSQL, defined by 21 Prisma models in
\texttt{backend/prisma/schema.prisma}. Prisma acts as the
abstraction layer: every query is written against a typed API such as
\texttt{findMany}, \texttt{upsert} or \texttt{createMany} and compiled
into parameterised SQL, which eliminates the possibility of injecting
raw SQL from application code.
}

\begin{table}[t]
\caption{\bi{21 model Prisma cùng số trường vô hướng và quan hệ.}{The 21
Prisma models with their scalar and relation field counts.}}
\label{tab:erd}
\centering
\scriptsize
\fitcol{% Auto-generated from backend/prisma/schema.prisma - do not edit by hand.
\begin{tabular}{@{}llrr@{}}
\toprule
\textbf{\bi{Model}{Model}} & \textbf{\bi{Vai trò nghiệp vụ}{Business role}} & \textbf{\bi{Scalar}{Scalar}} & \textbf{\bi{Quan hệ}{Relation}} \\
\midrule
\texttt{User} & \bi{Tài khoản người dùng + RBAC}{User account + RBAC} & 14 & 13 \\
\texttt{RefreshToken} & \bi{Xoay vòng phiên đăng nhập}{Rotating login sessions} & 6 & 1 \\
\texttt{PasswordResetToken} & \bi{Đặt lại mật khẩu (TTL 1h)}{Password reset (1h TTL)} & 7 & 1 \\
\texttt{ProductionTag} & \bi{Thẻ nội dung theo chương trình}{Content tag per production} & 5 & 0 \\
\texttt{LiveChannel} & \bi{Kênh phát sóng trực tiếp}{Broadcast channel} & 29 & 5 \\
\texttt{LiveEvent} & \bi{Sự kiện phát sóng / chương trình}{Live event / programme} & 32 & 7 \\
\texttt{Recording} & \bi{Bản ghi đã phát}{Published recording} & 42 & 5 \\
\texttt{Podcast} & \bi{Chuyên mục podcast}{Podcast collection} & 12 & 3 \\
\texttt{PodcastEpisode} & \bi{Tập podcast}{Podcast episode} & 12 & 1 \\
\texttt{Follow} & \bi{Theo dõi kênh}{Channel follow} & 4 & 2 \\
\texttt{Comment} & \bi{Bình luận}{Comment} & 11 & 6 \\
\texttt{CommentLike} & \bi{Like bình luận}{Comment like} & 4 & 2 \\
\texttt{Reaction} & \bi{Cảm xúc realtime}{Realtime reaction} & 6 & 3 \\
\texttt{WatchHistory} & \bi{Lịch sử xem}{Watch history} & 8 & 2 \\
\texttt{Notification} & \bi{Thông báo trong ứng dụng}{In-app notification} & 13 & 2 \\
\texttt{AuditLog} & \bi{Nhật ký kiểm toán (gRPC)}{Audit log (gRPC)} & 11 & 0 \\
\texttt{Watchlist} & \bi{Danh sách theo dõi}{Watchlist} & 7 & 1 \\
\texttt{AiCuratorReport} & \bi{Báo cáo AI Curator}{AI Curator report} & 15 & 1 \\
\texttt{DeviceToken} & \bi{Token thiết bị FCM}{FCM device token} & 8 & 1 \\
\texttt{NewsletterSubscriber} & \bi{Người đăng ký newsletter}{Newsletter subscriber} & 14 & 1 \\
\texttt{StreamDonation} & \bi{Quyên góp trong livestream}{Livestream donation} & 14 & 3 \\
\bottomrule
\end{tabular}
}
\end{table}

\subsection{\bi{Phân nhóm miền dữ liệu}{Data-Domain Grouping}}

\bi{
Các model được tổ chức theo năm nhóm nghiệp vụ:
}{
The models are organised into five business groups:
}

\begin{itemize}[leftmargin=*,itemsep=1pt]
  \item \bi{\textbf{Xác thực và danh tính.} \texttt{User} lưu thông tin
  tài khoản cùng vai trò, \texttt{RefreshToken} lưu phiên đăng nhập
  xoay vòng, \texttt{PasswordResetToken} lưu mã đặt lại mật khẩu với
  thời hạn một giờ.}{
  \textbf{Authentication and identity.} \texttt{User} stores account
  details together with the role, \texttt{RefreshToken} stores
  rotating login sessions, and \texttt{PasswordResetToken} stores
  password-reset codes with a one-hour lifetime.}

  \item \bi{\textbf{Nội dung phát sóng.} \texttt{LiveChannel} là kênh,
  \texttt{LiveEvent} là chương trình trong lịch, \texttt{Recording}
  là bản ghi đã phát, \texttt{ProductionTag} là thẻ nội dung gắn với
  chương trình.}{
  \textbf{Broadcast content.} \texttt{LiveChannel} is the channel,
  \texttt{LiveEvent} is a scheduled programme, \texttt{Recording} is
  a published recording and \texttt{ProductionTag} is a content tag
  attached to a production.}

  \item \bi{\textbf{Tương tác người dùng.} \texttt{Comment},
  \texttt{CommentLike}, \texttt{Reaction}, \texttt{Follow},
  \texttt{WatchHistory}, \texttt{Watchlist} và
  \texttt{StreamDonation} mô tả toàn bộ hành vi của người xem.}{
  \textbf{User interaction.} \texttt{Comment}, \texttt{CommentLike},
  \texttt{Reaction}, \texttt{Follow}, \texttt{WatchHistory},
  \texttt{Watchlist} and \texttt{StreamDonation} describe the full
  spectrum of viewer behaviour.}

  \item \bi{\textbf{Vận hành và kiểm toán.} \texttt{AuditLog} nhận
  bản ghi từ microservice gRPC, \texttt{Notification} lưu thông báo
  trong ứng dụng, \texttt{DeviceToken} lưu token FCM,
  \texttt{NewsletterSubscriber} lưu người đăng ký bản tin.}{
  \textbf{Operations and auditing.} \texttt{AuditLog} receives records
  from the gRPC microservice, \texttt{Notification} stores in-app
  notifications, \texttt{DeviceToken} stores FCM tokens and
  \texttt{NewsletterSubscriber} stores newsletter subscribers.}

  \item \bi{\textbf{Trợ lý AI.} \texttt{AiCuratorReport} lưu kết quả
  biên tập AI kèm dấu vết phiên bản prompt và mô hình.}{
  \textbf{AI assistance.} \texttt{AiCuratorReport} stores AI curation
  output together with its prompt and model version provenance.}
\end{itemize}

\subsection{\bi{Quan hệ và ràng buộc}{Relationships and Constraints}}

\bi{
Quan hệ được định nghĩa bằng khóa ngoại chuẩn của Prisma, và mỗi quan
hệ đều có chỉ mục phù hợp với truy vấn hot-path. Ba ràng buộc đáng
chú ý:
}{
Relationships are expressed with standard Prisma foreign keys, and
each carries an index matched to its hot-path query. Three constraints
deserve attention.
}

\begin{enumerate}[leftmargin=*,itemsep=1pt]
  \item \bi{\textbf{Xoá người dùng không phá vỡ nhật ký kiểm toán.}
  \texttt{RefreshToken.user} có quan hệ \texttt{OnDelete: SetNull},
  nên xoá tài khoản chỉ vô hiệu hoá phiên mà vẫn giữ được dấu vết
  kiểm toán.}{
  \textbf{User deletion does not destroy the audit trail.}
  \texttt{RefreshToken.user} uses \texttt{OnDelete: SetNull}, so
  deleting an account invalidates the sessions while preserving the
  audit trail.}

  \item \bi{\textbf{Khoá duy nhất trên thẻ nội dung.} \texttt{ProductionTag}
  có khoá duy nhất trên cặp chương trình và tên thẻ, nên cùng một thẻ
  không thể bị gán hai lần.}{
  \textbf{A unique key on content tags.} \texttt{ProductionTag} carries
  a unique constraint on the programme-and-tag pair, so the same tag
  cannot be attached twice.}

  \item \bi{\textbf{Chỉ mục theo trục thời gian.} Các truy vấn EPG
  luôn lọc theo \texttt{scheduledAt} và \texttt{channelId}, nên cột
  này được chỉ mục; truy vấn lịch sử xem lọc theo
  \texttt{userId} và \texttt{watchedAt}.}{
  \textbf{Time-axis indexes.} EPG queries always filter by
  \texttt{scheduledAt} and \texttt{channelId}, so those columns are
  indexed; watch-history queries filter by \texttt{userId} and
  \texttt{watchedAt}.}
\end{enumerate}

\subsection{\bi{Chuẩn hoá và quyết định thiết kế}{Normalisation and Design Decisions}}

\bi{
Mô hình được chuẩn hoá ở mức 3NF: các thuộc tính chỉ phụ thuộc khóa
chính, không có phụ thuộc hàm không chuyển tiếp giữa các thuộc tính
không khóa. Hai quyết định thiết kế đáng bàn:
}{
The model is normalised to third normal form: attributes depend only on
the primary key and no non-key attribute depends transitively on
another. Two design decisions deserve discussion.
}

\begin{itemize}[leftmargin=*,itemsep=1pt]
  \item \bi{\textbf{Chuẩn hoá cấu trúc EPG.} Thay vì lưu lưới 24
  giờ đã điền, hệ thống lưu các sự kiện thật cùng thẻ nội dung, và
  lưới được tính lại khi truy vấn. Cách này giữ dữ liệu gốc trung
  thực và cho phép lịch thay đổi mà không cần di chuyển dữ liệu, đổi
  lại tốn công tính toán tại thời điểm đọc --- một đánh đổi đúng đắn
  khi phần lớn thao tác ghi là chương trình biên tập và phần lớn
  thao tác đọc là xem lịch.}{
  \textbf{EPG structure normalisation.} Rather than storing the
  filled 24-hour grid, the system stores real events with content tags
  and recomputes the grid on demand. This keeps the source data
  truthful and lets schedules change without data migration, at the
  cost of computation at read time --- a sound trade-off when most
  writes are editorial and most reads are schedule viewing.}

  \item \bi{\textbf{Không lưu sẵn báo cáo AI trong JSON blob.}
  \texttt{AiCuratorReport} có các cột riêng cho từng trường quan
  trọng (điểm tin cậy, phân loại phù hợp phát sóng, cờ tuân thủ)
  để có thể lọc và thống kê bằng SQL, chỉ giữ phần phân tích tự do
  dạng văn bản trong trường riêng.}{
  \textbf{AI reports are not stored as a JSON blob.}
  \texttt{AiCuratorReport} exposes separate columns for each important
  field (confidence score, broadcast-suitability class, compliance
  flags) so they can be filtered and aggregated in SQL, keeping only
  the free-form analysis in a dedicated text column.}
\end{itemize}

\subsection{\bi{Quản lý phiên bản lược đồ}{Schema Versioning}}

\bi{
Mọi thay đổi lược đồ đều được quản lý bằng \texttt{prisma migrate},
tạo ra các tệp SQL trong \texttt{backend/prisma/migrations} được đưa
vào kiểm soát phiên bản. CI chạy \texttt{prisma migrate deploy} trên
một dịch vụ PostgreSQL 15 tạm trước khi biên dịch ứng dụng, nhờ vậy
một lược đồ không tương thích sẽ bị phát hiện trước khi triển khai
sản xuất.
}{
Every schema change is managed through \texttt{prisma migrate}, which
generates SQL files under \texttt{backend/prisma/migrations} that are
kept under version control. CI runs \texttt{prisma migrate deploy}
against a temporary PostgreSQL 15 service before building the
application, so an incompatible schema is detected before any
production deployment.
}
